\subsection{The completed thimbles and the marked fibre}\label{sec:completedthimbles}

We calculate the morphisms in the completion stopped at the core of $F_0$. Here the core of a Liouville domain is the set of points which do not escape under positive Liouville flow. The core of its boundary hypersurface defines a stop in the contact boundary at infinity. There are many equivalent descriptions of a stopped category. For the present calculation we need one that preserves the individual disks: the order of their ends controls which Floer complex computes their morphisms.

\begin{proposition}\label{prop:diskhoms}
At a fixed Morse parameter one can equip the common-path thimbles with gradings and relative Pin structures so that their isomorphism classes are permuted by $G$. In the category stopped at the core of $F_0$ each disk has endomorphism cohomology $\C$ in degree zero. With the outgoing paths in increasing terminal-angle order, all backward morphism groups vanish. Distinct disks at a common critical value are mutually orthogonal in every degree.
\end{proposition}

We give the changes at the two ends of the fibration before making the Floer calculation. The first puts a radial product at infinity without changing the marked page. The second places all the thimbles in the half-plane chart used by Ganatra, Pardon and Shende.

\begin{proof}
We may standardize the Morse charts before choosing the paths in Lemma~\ref{lem:samedisks}. Here are details ensuring that this step respects the stabilizers. Average a local holomorphic chart, composed in each summand with the inverse tangent action, to linearize the stabilizer. The quadratic jet is a nondegenerate invariant complex quadratic form. In an invariant Hermitian metric its associated antilinear operator has antiunitary polar factor a real structure: in a unitary Takagi basis the operator is complex conjugation times a positive diagonal matrix. On the fixed real subspace the quadratic form is positive definite. Real orthonormal coordinates for that form give $q=\sum z_i^2$ with real orthogonal stabilizer action. To remove the terms $h$ of order at least three, solve $d(q+th)(v_t)=-h$. The gradient entries of $q+th$ are coordinates near zero, uniformly for $0\le t\le1$. In them $h$ is a sum of the entries times functions of order at least two, by integration along radii. Average the resulting holomorphic field and integrate on a smaller ball.

After removing the real pluriharmonic part of its quadratic Taylor polynomial, an invariant potential $P$ for the local Kähler form has positive Hermitian quadratic term. Let $P_0$ give the standard form, proportional to $|z|^2$. Replace $P$ by $P+t\chi(P_0-P)$ in each of the disjoint charts. Choose the radial cutoff $\chi$ equal to one near zero and zero outside the chart, with $|r\chi'|+|r^2\chi''|$ small by taking a long transition in $\log r$. Since $P_0-P=O(r^2)$ and $d(P_0-P)=O(r)$, the cutoff errors in the Levi form are smaller than the uniform positive lower bound of the convex combination. This gives a deformation through invariant primitives by adding $d^c(t\chi(P_0-P))$ to the fibrewise construction before the base-area addition. In the added base term use $\kappa b(|w|)\alpha_0$ with $b$ nondecreasing, zero on a disk containing the Morse charts and one on an outer annulus. This changes the primitive through nonnegative additions of base area, preserves the outward estimates and makes the \emph{total} form standard at each centre. On the product horizontal collar, $\alpha_R-\alpha_0=df$ with $f(R)=0$. Subtract $d(\nu f(\pi))$ using a fibre cutoff $\nu$ equal to one near the face and supported away from the disks. There the resulting product primitive is $\lambda_F+(1+\kappa b)\alpha_0$. The marked fibre retains its restricted primitive, and the compact vertical-face errors are bounded before increasing $\kappa$. In the separated fibration of Lemma~\ref{lem:samedisks} we now subtract $(\pi')^*((1+\kappa b)\alpha_0)$ to obtain the fibrewise primitive. This gives $\lambda_F$ on the horizontal collar. Near each Morse centre $b=0$ and the subtracted area is of type $(1,1)$ and vanishes at the centre, so the fibrewise form is Kähler on a smaller chart. On all radial ends $b=1$. The radial change of base coordinate $w\mapsto\sqrt{1+\kappa b(|w|)}\,w$ makes the exterior base primitive standard. We keep the old coordinate until the next calculation is finished.

On a short annulus $R_0\le r\le R$ trivialize in the radial direction by parallel transport for a separated presentation $\pi'$, with total primitive $\Lambda$. This follows each actual disk in its radial collar. Set $K=1+\kappa$ and $\Theta=\Lambda-K(\pi')^*\alpha_0$. Let $\zeta(r)$ be smooth and nondecreasing, equal to $r$ near $R_0$ and to $R$ near $R$. Write $j_s$ for the substitution of $\zeta_s(r)=(1-s)r+s\zeta(r)$ for $r$ in this trivialization. On the annulus replace $\Lambda$ by
\[
                         j_s^*\Theta+K(\pi')^*\alpha_0.
\]
The fibre form is constant in these radial coordinates and $d\Theta$ pairs the radial vector with each vertical vector to zero. If $a$ denotes the original horizontal area coefficient relative to $r\,dr\wedge d\theta$, the new coefficient is
\[
              K+\zeta_s'(r)\frac{\zeta_s(r)}{r}
                              \big(a(\zeta_s(r))-K\big).
\]
For the unsplit map $a-K$ is uniformly bounded independently of $K$; adding base area preserves the connection. Choose the annulus and cutoff first, then $K$ large enough for positivity there. For this fixed $K$, sufficiently small splitting times preserve strict positivity by smooth dependence of radial transport on the compact annulus. There are finitely many refinements. The tangential primitive at $r=R$ and the horizontal product collar stay fixed, including the corner rounding. The radial substitution takes each terminal collar into itself, where the primitive vanishes on the disk. Thus the disks remain Lagrangian with Legendrian ends. Identify the completed collars by boundary point and Liouville time. Exact Moser transport in the compact part then gives an exact isotopy of the disks with fixed cylindrical ends avoiding the stop, hence identical objects up to isomorphism \cite[Remark~2.30]{GPS}.

At $s=1$ continue the primitive as $\beta+p\,d\theta$, $p=Kr^2/2$, where $\beta$ is pulled back from the bundle on the outer circle. Attach the product fibre end as well. On the universal cover of the circle use the kernel of $d\beta$ for angular transport, starting at the reference fibre at $\theta=0$. Thus $d\beta=\omega_F$, and
\[
       \beta=\lambda_F+dG,\qquad
       G(\theta,x)=\int_0^\theta\beta(\partial_\vartheta)(\vartheta,x)\,d\vartheta .
\]
This follows by differentiating along the fibre, using $\beta|_{TF}=\lambda_F$ at $\theta=0$. The integral fixes the constant on every fibre component and is zero on the product fibre end. If $\mu$ is the monodromy, $k(x)=G(2\pi,x)$ satisfies $\mu^*\lambda_F-\lambda_F=dk$ and
$G(\theta+2\pi,x)=k(x)+G(\theta,\mu x)$. To flatten the primitive away from a short angular interval, choose $G_0=\rho(\theta)k(x)$ on one strip, with $\rho$ zero near zero and one near $2\pi$, extending by this same relation. Smoothness at the seam uses constancy of $\rho$ near both ends. Then $D=G_0-G$ descends to the circle bundle and has compact fibre support.

Change the primitive by $d(\chi_P(p)D)$, with $\chi_P$ zero up to a large $P$ and one beyond $2P$. The two-form is fixed. The tails of the geometric disks are products $V_i\times\{\theta_i\}\times[P,\infty)$ with $V_i$ compact exact vanishing spheres. To keep a disk cylindrical through the change, subtract the differential of an extension of the \emph{function} $\chi_P(p)D(\theta_i,x)$ on its tail, cut off in the fibre and in angle away from the marked ray. This preserves the primitive on the disk, also on the transition in $p$. For a zero-dimensional vanishing sphere the two restriction constants may differ. Outward Liouville flow identifies the resulting cylindrical representative with the usual cylindrization for the new uncorrected primitive \cite[\S7.2, (7.7)--(7.9)]{GPS}.

To check the stop under this identification, write the primitive as $p\,d\theta+\lambda_F+dg$ in angular transport coordinates. Its Liouville field satisfies $\dot p=p+\partial_\theta g$ and $\dot\theta=-\partial_p g$. All angular derivatives of the corrections are bounded on a fundamental strip. Take $P$ large so $\dot p\ge p/2$ in the region of change. Before it, outwardness at the original base boundary persists because $\partial_\theta G$ is independent of $p$. The fibre end is a product. These estimates give complete outward flows from a common compact carrier, whose conjugacy identifies the primitives strictly. On the reference ray $G(0,x)=G_0(0,x)=0$, all disk cutoffs vanish, $\dot\theta=0$, and the vertical field is exactly $Z_F$. Its core is invariant in both directions. The stop is therefore the same core of $F_0$. Applying the same conjugacy throughout the interpolation transports the cylindrical disks by an exact cylindrical isotopy. The calculation below can be made in the resulting flat model.

\medskip
\noindent\emph{A half-plane containing the outgoing paths.}
Write $0<\theta_1<\cdots<\theta_\mu$ for the small terminal angles. Choose $0<a<\theta_1$ and put the transition of $\rho$ strictly inside $(0,a)$. In the angular trivialization based backwards at $2\pi$, the primitive at large base radius is exactly $\lambda_F+p\,d\theta$ on a neighbourhood of $[a,2\pi]$. Indeed the fibre primitive after the transition on the original strip is $\mu^*\lambda_F$. The backwards trivialization removes this pullback and crosses the seam where the new potential is zero. It is important to put the positive-real stop at the \emph{upper} end of the right half-plane in \cite[\S8.6]{GPS}.

Let $f$ be an increasing circle diffeomorphism carrying $[-\pi/2,\pi/2]$ to $[a,2\pi]$ in this lift. The map $(p,\theta)\mapsto(p/f'(\theta),f(\theta))$ preserves $p\,d\theta$ and takes rays to rays. To place also the finite part of the half-plane in the exterior chart, first apply a compactly supported Hamiltonian translation. In Cartesian coordinates take $Cy\,\chi(y/T)\psi(x/T)$ with the sign translating by $+C$ along the real axis, with both cutoffs one on $[-1,1]$ and zero outside $[-2,2]$. Choose $C$ large compared with the inner radius of the chart times $\sqrt{\max f'}$. The time-one displacement is bounded by a constant times $C$, independently of $T$. For $T$ sufficiently large the half-disk of radius $T/2$ is translated by $+C$. Points outside it stay at large radius and change angle by $O(C/T)$, within the available angular margin. Composing with the polar map gives a proper embedding $j$ of the closed right half-plane into the chart. Its upper boundary end maps to the marked ray and all the thimble tails lie on the images of interior rays. The primitive difference is $j^*\alpha_0-\alpha_0=dh$ for a compactly supported function $h$, by compact Hamiltonian isotopy.

On a noncompact fibre a correction compactly supported in the base alone would not give the required Liouville embedding at total-space infinity. Extend $h\circ j^{-1}$ from the closed image to a smooth compactly supported function $h'$ of the target base, using a collar of its boundary. In the product fibre end subtract $d(\nu h'(\pi'))$, where $\nu$ is zero on a large compact fibre set and one past a larger level. The resulting difference on the product half-plane chart is the differential of $(1-\nu)h$, compactly supported in the total space. In fibre-end coordinates $I\alpha_F$ the added fibre covector is a multiple of $dI$, so evaluating the Liouville equation on the fibre Reeb field still gives $\dot I=I$. The change vanishes near a compact carrier with base radius above the support of $h'$ and fibre radius below the support of $\nu$. Fibre radial growth and the unchanged base growth beyond that support identify completions by their outward flows. The fibre cutoff can be placed beyond all the disks over the compact base support. Around their tails and the marked core the previous primitives and cylindrical identifications persist. For a finite fibre the product correction already has compact support.

\medskip
\noindent\emph{Floer cohomology.}
Each disk now has, outside a compact set, a compact exact fibre slice times a ray in this half-plane. Thus it is a thimble-like object of \cite[Definition~8.19]{GPS}. By full faithfulness in \cite[Proposition~8.20]{GPS}, its morphisms in the stopped category can be computed by moving the source in increasing terminal angle above the target and using ordinary Floer cohomology there. For $\theta_i>\theta_j$ the source is already above the target. The two disks lie over disjoint paths, and a sufficiently small positive pushoff at the source end still misses the target. The complex has no generators.

For a self-morphism move the path slightly to its left, rotating it by a small positive angle $\delta$ at its Morse centre and increasing its terminal angle. A normal strip along the path gives such a pushoff meeting the old path only at the centre. Transporting the vanishing sphere through these paths gives a positive thimble-like isotopy; its compact slice may vary from one path to the next. At the centre the two tangent planes in the standard quadratic chart are $\mathbb R^n$ and $e^{i\delta/2}\mathbb R^n$, after fixing the initial phase. The whole Floer complex therefore has one generator and zero differential. Its endomorphism cohomology is a nonzero unital graded algebra, so its one dimension lies in degree zero. Lemma~\ref{lem:samedisks} allows both orders at each common value. The two backward vanishings obtained this way concern the same original disks.

Finally start on the logarithmic torus with the square of its logarithmic volume form. A fan automorphism changes the volume by a determinant $\pm1$, so this square is invariant. Extend it along the invariant deformations used above: choose invariant compatible metrics and transport the squared determinant line along the interval, then extend on Liouville collars. For relative Pin data take the product real bundle with the fan representation on its fibre, with background its second Stiefel--Whitney class. Each contractible disk has a unique relative Pin class up to isomorphism. Its circle-valued squared phase has a lift, and a stabilizer element changes a lift by a constant integer. Finite order forces that integer to be zero. Taking one lift per orbit and transporting gives the asserted action on the isomorphism classes of the decorated disks in the initial presentation. Only the auxiliary calculations after separation required a choice of order within a common value.
\end{proof}
